% !TeX root = ../../Book3.tex
% 纲要标题：### E.4 线性微分理论与相关专题

\section{线性微分理论与相关专题}
\label{b3:sec:E04}

线性微分方程把求导写成有限维空间上的加法算子，但该算子对标量满足带导子的乘法法则。解矩阵与常数域因此共同决定对称群；加入参数后，还要记录额外导子之间的相容关系。

% 纲要细目（与计划纲要同步）：
% * 线性系统、微分模及 Picard–Vessiot 理论转第五卷附录 E，不要求先掌握全部特征集算法
% * 微分闭域、参数化理论、微分代数群及差分代数只作后续地图；参数化情形的常数域假设另列
% * 完整微分代数几何、Kolchin 维数多项式、算法复杂度及微分模型论不纳入本套书的共同证明链
%
% #### E.4.1 线性系统、微分模与 Picard–Vessiot 理论
% #### E.4.2 微分闭域、参数化理论、微分代数群与差分代数
% #### E.4.3 微分代数几何、维数多项式与模型论导读


% 正文已按纲要写入；有限性与算法长证明给出概要及出处。

\subsection{线性系统、微分模与 Picard–Vessiot 理论}
\label{b3:subsec:E04:01}

\begin{definition}\label{b3:def:differential-module}
设 \((K,\delta)\) 为特征零微分域，\(M\) 为有限维 \(K\) 向量空间。若加法映射 \(\nabla:M\to M\) 满足 \(\nabla(am)=\delta(a)m+a\nabla(m)\)，则 \((M,\nabla)\) 称为\term{微分模}[differential module]。满足 \(\nabla(m)=0\) 的向量称为水平向量。
\end{definition}

选定基后可写 \(\nabla(v)=\delta(v)-Av\)，于是水平向量满足线性系统 \(Y'=AY\)。若令 \(Y=PZ\)，\(P\in\operatorname{GL}_n(K)\)，则新系统为 \(Z'=(P^{-1}AP-P^{-1}P')Z\)。额外的 \(P'\) 项来自基随导子变化，故不能照普通矩阵相似变换处理。

\begin{definition}\label{b3:def:picard-vessiot-extension}
设 \((K,\delta)\) 为特征零微分域，\(C=\ker\delta\) 为常数域，\(n\ge1\)，\(A\in M_n(K)\)。若微分域扩张 \(L/K\) 含有矩阵 \(Y\in\operatorname{GL}_n(L)\)，满足 \(Y'=AY\)，且 \(L\) 由 \(Y\) 的各项与 \(\det(Y)^{-1}\) 在 \(K\) 上生成为域，并有 \(C_L=C\)，则称 \(L/K\) 为该系统的\term{Picard--Vessiot 扩张}[Picard--Vessiot extension]。与导子交换且固定 \(K\) 的 \(L\) 的域自同构组成其微分 Galois 群。
\end{definition}

\begin{proposition}\label{b3:prop:pv-constant-matrices}
设 \((K,\delta)\) 为特征零微分域，\(C=\ker\delta\)，\(A\in M_n(K)\)，\(L/K\) 为方程 \(Y'=AY\) 的 Picard--Vessiot 扩张，且选定使 \(L=K(Y_{ij},\det(Y)^{-1})\) 的解矩阵 \(Y\in\operatorname{GL}_n(L)\)。每个与导子交换且固定 \(K\) 的自同构 \(\sigma:L\to L\) 唯一确定 \(C_\sigma\in\operatorname{GL}_n(C)\)，使 \(\sigma(Y)=YC_\sigma\)。映射 \(\sigma\mapsto C_\sigma\) 是单射群同态。
\end{proposition}
\begin{proof}
\(\sigma(Y)'=A\sigma(Y)\)，而 \((Y^{-1})'=-Y^{-1}A\)，故 \((Y^{-1}\sigma(Y))'=0\)。无新常数条件使矩阵各项属于 \(C\)。它可逆并唯一。复合时常数矩阵被自同构固定，于是 \(C_{\sigma\tau}=C_\sigma C_\tau\)。若 \(C_\sigma=I\)，则 \(\sigma\) 固定域生成元，故为恒等。
\end{proof}

一般存在唯一性、群的代数性及 Galois 对应需要额外证明，本附录不把以上矩阵嵌入当作那些定理。第一卷附录 D 的基本解与微分自同构给出初步算例；第五卷相应专题将从表示与线性系统继续讨论。

\subsection{微分闭域、参数化理论、微分代数群与差分代数}
\label{b3:subsec:E04:02}

微分方程未必在给定微分域中有解，这促使人们扩大“代数闭”的观念。设 \((K,\delta)\) 为只指定一个导子的特征零微分域。若任意以 \(K\) 为系数、在某个微分扩域可解的有限微分多项式方程与不等式系统，在 \(K\) 中已有解，则称 \((K,\delta)\) 为\term{微分闭域}[differentially closed field]。这里不等式指 \(g\ne0\)，并非序关系。只有一个导子的理论与多个交换导子的理论必须分别注明；它们都不等同于底层域代数闭。

\ProseParagraphBreak
若方程还依赖参数，可以同时指定主变量导子 \(\delta\) 和参数导子 \(\partial\)，要求二者交换。\(\delta\) 的常数域对 \(\partial\) 稳定，因为 \(\delta a=0\) 蕴含 \(\delta(\partial a)=\partial(\delta a)=0\)；它自身因而带有参数导子。研究与二者同时交换的自同构时，常数矩阵的系数还会受到 \(\partial\) 方程约束，这就是参数化微分 Galois 理论出现微分代数群的缘由。所谓\term{微分代数群}[differential algebraic group]，在矩阵情形可理解为由微分多项式方程定义、且对乘法和取逆封闭的矩阵集合。例如常数非零元组成 \(\{a:a\ne0,\delta a=0\}\)，它是乘法群的微分代数子群。这段说明没有提出或使用参数化扩张的存在性定理。

\ProseParagraphBreak
\term{差分代数}[difference algebra]则指定保持乘法的自同态 \(\sigma\)，如 \(f(t)\mapsto f(t+1)\)；它研究 \(y(t+1)\) 与 \(y(t)\) 的关系。导子的 Leibniz 法则与自同态的乘法法则不同，因而差分多项式与微分多项式的理想条件并不相同。本节只说明对象与问题，未建立这些专题的存在定理或完整 Galois 对应。

\subsection{微分代数几何、维数多项式与模型论导读}
\label{b3:subsec:E04:03}

微分代数几何用有限阶截断，把无穷多个导数暂时收集在有限生成的代数中。设 \(L=K\langle a_1,\ldots,a_n\rangle\)，定义
\[
h_a(s)=\operatorname{trdeg}_K K(a_i,\delta a_i,\ldots,\delta^sa_i:1\le i\le n).
\]
它统计到第 \(s\) 阶为止还保留多少个普通代数自由参数。若 \(a\) 是一个微分超越元，则 \(h_a(s)=s+1\)；若 \(a'=a\) 且 \(a\) 在 \(K\) 上普通超越，则所有高阶导数等于 \(a\)，故 \(h_a(s)=1\)。两种情形都不能只用“方程有几个未知函数”判断。

\ProseParagraphBreak
若一个数值多项式在充分大的整数 \(s\) 上等于 \(h_a(s)\)，就称它为这一生成元组的\term{微分维数多项式}[differential dimension polynomial]。上面两个例子分别得到 \(s+1\) 与 \(1\)；多个交换导子时则按总微分阶截断。一般存在性定理及其与排序、特征集的关系属于本附录未建立的维数理论。有效计算所涉及的问题可读\cite[Section 1]{Hubert2003DifferentialSystems}。截断过程记录的是给定生成元的导数关系，不能把任意有限阶的解都认作可延拓为全部阶的解。

\ProseParagraphBreak
模型论把这些问题写成微分域上的公式：方程表达等式，不等条件表达可逆性，量词表达存在新的解或参数。微分闭域提供研究解集的一种环境，但它不自动指定解析函数，也没有给出收敛区域。读到这里所应掌握的，是如何从方程得到微分理想、如何在消去中保留非零条件，以及线性系统怎样产生常数矩阵作用；更一般的几何与模型论需要另作系统学习。
